\section{Results}
\label{sec:results}

We evaluate whether the logit-space disagreement between a correctness probe
and an ambiguity/expert-reliability probe improves selective prediction. The
main results are organized around three questions: whether the gap improves
risk--coverage behavior across models and datasets, whether the diagnostic
regime collapses as predicted, and whether the expert-reliability signal is
load-bearing rather than an incidental label source.

\subsection{Headline Risk--Coverage Matrix}
\label{sec:results-headline}

Table~\ref{tab:headline-matrix} reports AURC and Coverage@Acc across
model--dataset pairs and scoring methods. Lower AURC is better; higher
Coverage@Acc is better. We report mean and standard deviation over five seeds.
Bold marks the best method per row. Asterisks mark paired-permutation
significance against the next-best method at $p < 0.01$.

\begin{table*}[t]
\centering
\caption{\textbf{Headline selective-prediction results.} Rows are organized
with showcase datasets first, followed by diagnostic datasets. Placeholder
columns should be filled from \texttt{outputs/aggregated\_results.parquet}.}
\label{tab:headline-matrix}
\begin{tabular}{llrrrrr}
\toprule
Model & Dataset & Logprob & Seq. ent. & Probe-only & Sem. ent. & Gap \\
\midrule
% Llama-3.1-8B & ChaosNLI & -- & -- & -- & -- & -- \\
% Llama-3.1-8B & Pavlick-NLI & -- & -- & -- & -- & -- \\
% Llama-3.1-8B & AmbigQA-k$\geq$2 & -- & -- & -- & -- & -- \\
\midrule
% Llama-3.1-8B & TruthfulQA & -- & -- & -- & -- & -- \\
% Llama-3.1-8B & HaluEval & -- & -- & -- & -- & -- \\
\bottomrule
\end{tabular}
\end{table*}

\subsection{Risk--Coverage Curves}
\label{sec:results-rc-curve}

Figure~\ref{fig:chaosnli-rc} shows the full risk--coverage curve on the
canonical showcase setting, ChaosNLI with the headline model. The shaded band
is the seed-aggregated 95\% confidence interval. Vertical markers indicate the
coverage achieved at Accuracy@0.90 thresholds from
Table~\ref{tab:headline-matrix}.

\begin{figure}[t]
\centering
% \includegraphics[width=\linewidth]{figures/chaosnli_rc_curve.pdf}
\caption{\textbf{Risk--coverage curve on ChaosNLI.} The gap should dominate
single-signal baselines across the coverage range, not only in a narrow band.}
\label{fig:chaosnli-rc}
\end{figure}

\subsection{Diagnostic Regime Collapse}
\label{sec:results-diagnostic}

Table~\ref{tab:diagnostic-collapse} reports the AURC delta between the gap and
the single correctness probe, together with the effective dimensionality of the
$(p_{\mathrm{pred}}, p_{\mathrm{amb}})$ cloud. The SCOD prediction is that
datasets with effectively one-dimensional probe geometry should show little or
no gain from the gap, while showcase datasets with genuine two-dimensional
structure should benefit.

\begin{table}[t]
\centering
\caption{\textbf{Diagnostic-regime collapse.} Positive $\Delta$ means the gap
improves over the single correctness probe. Effective dimensionality comes from
\texttt{scripts/scod\_diagnostics.py}.}
\label{tab:diagnostic-collapse}
\begin{tabular}{llrr}
\toprule
Model & Dataset & $\Delta$ AURC & Eff. dim. \\
\midrule
% -- & TruthfulQA & -- & -- \\
% -- & HaluEval & -- & -- \\
\bottomrule
\end{tabular}
\end{table}

\subsection{\texorpdfstring{$y_{\mathrm{expert}}$}{y expert} Semantic Audit}
\label{sec:results-yexpert-audit}

Table~\ref{tab:yexpert-audit} tests whether the expert-reliability labels are
semantically load-bearing. The native signal should improve over the
single-probe baseline; shuffling within correctness strata and replacing the
signal with a constant should collapse toward the single-probe baseline; the
inverse-oracle variant should dominate as a sanity check.

\begin{table*}[t]
\centering
\caption{\textbf{Semantic audit of the expert-reliability signal.} Lower AURC
is better. This is the main-text defense that the second head is not merely an
incidental label source.}
\label{tab:yexpert-audit}
\begin{tabular}{llrrrr}
\toprule
Model & Dataset & Native & Shuffled & Constant & Inverse oracle \\
\midrule
% -- & ChaosNLI & -- & -- & -- & -- \\
% -- & Pavlick-NLI & -- & -- & -- & -- \\
\bottomrule
\end{tabular}
\end{table*}

\subsection{SCOD-Optimal Selector Check}
\label{sec:results-scod-selector}

We compare the gap against calibrated SCOD-additive scoring and a learned 2D
selector trained on $(\mathrm{logit}\,p_{\mathrm{pred}},
\mathrm{logit}\,p_{\mathrm{amb}})$. This experiment is decision-revealing: if
the learned 2D selector substantially outperforms the gap, the headline method
should switch before submission.

\begin{table}[t]
\centering
\caption{\textbf{SCOD selector comparison.} This table can remain appendix-only
unless the learned 2D selector changes the headline-method decision.}
\label{tab:scod-selector}
\begin{tabular}{llrrrr}
\toprule
Model & Dataset & Gap & Gap-cal. & SCOD-add. & Learned 2D \\
\midrule
% -- & ChaosNLI & -- & -- & -- & -- \\
\bottomrule
\end{tabular}
\end{table}

\subsection{Conformal Validation}
\label{sec:results-conformal}

Table~\ref{tab:conformal-validation} validates that the conformal thresholding
step achieves the intended selective-accuracy target within finite-sample
tolerance. The operational quantity is committed coverage at each target.

\begin{table}[t]
\centering
\caption{\textbf{Conformal validation.} Fill from the conformal validation
sweep at $\alpha \in \{0.10, 0.20\}$.}
\label{tab:conformal-validation}
\begin{tabular}{llrrr}
\toprule
Model & Dataset & Target acc. & Realized acc. & Coverage \\
\midrule
% -- & -- & -- & -- & -- \\
\bottomrule
\end{tabular}
\end{table}

